% Compile on Overleaf with pdfLaTeX
\documentclass[12pt,a4paper]{article}
\usepackage[utf8]{inputenc}
\usepackage[T5]{fontenc}
\usepackage[vietnamese]{babel}
\usepackage{geometry}
\geometry{top=2.5cm,bottom=2.5cm,left=3cm,right=2cm}
\usepackage{setspace}
\onehalfspacing
\usepackage{hyperref}
\usepackage{xurl}
\hypersetup{colorlinks=true,linkcolor=black,urlcolor=blue,citecolor=black}
\setlength{\parindent}{1.25cm}
\setlength{\parskip}{0.2em}
\begin{document}
\begin{center}
{\Large\bfseries Cơ sở lý thuyết và công nghệ}
\end{center}

\section{Cơ sở lý thuyết}

Cơ sở lý thuyết của dự án tập trung vào dữ liệu điểm quan tâm (POI) trong du lịch số, quy trình xử lý dữ liệu nhiều nguồn và các nguyên tắc bảo đảm dữ liệu đủ tin cậy để phục vụ dashboard, API và các chức năng phân tích \cite{zhang2023}. Điểm chung của các nội dung này là dữ liệu POI không chỉ cần được thu thập, mà còn phải được chuẩn hóa, kiểm tra chất lượng, truy vết nguồn gốc và duy trì nhất quán trong suốt vòng đời xử lý \cite{ibmDataQuality}.\par

Phạm vi của dự án tập trung vào xây dựng nền tảng dữ liệu POI phục vụ dashboard vận hành và API dữ liệu, chưa đặt mục tiêu tối ưu thuật toán gợi ý điểm đến trong phiên bản hiện tại \cite{zhang2023}. Cách giới hạn phạm vi này giúp dự án ưu tiên chất lượng dữ liệu đầu vào, tính truy vết và khả năng phục vụ dữ liệu ổn định trước khi mở rộng sang các mô hình cá nhân hóa hoặc đề xuất nâng cao \cite{ibmDataQuality}.\par

\subsection{Dữ liệu điểm quan tâm và vấn đề tích hợp}

Dữ liệu điểm quan tâm mô tả các địa điểm có giá trị trong hành trình du lịch như nhà hàng, khách sạn, quán cà phê, điểm tham quan, công viên hoặc khu mua sắm \cite{zhang2023}. Trong hệ thống du lịch số, dữ liệu này là nền tảng cho tìm kiếm địa điểm, phân tích phân bố dịch vụ, dashboard quản trị và các chức năng gợi ý điểm đến \cite{zhang2023}.\par

Thách thức chính của dữ liệu POI là tính không đồng nhất giữa các nguồn \cite{yeboah2021}. Cùng một địa điểm có thể xuất hiện ở OpenStreetMap, dịch vụ thương mại hoặc nguồn nội bộ với tên gọi, tọa độ, địa chỉ, danh mục và thông tin liên hệ khác nhau \cite{yeboah2021}. Vì vậy, nền tảng dữ liệu du lịch cần xử lý các bước tích hợp, chuẩn hóa, loại bỏ trùng lặp, bổ sung thông tin thiếu và đánh giá độ tin cậy trước khi đưa dữ liệu vào lớp phục vụ \cite{ibmDataQuality}.\par

Trong bối cảnh dữ liệu du lịch Việt Nam, tên địa điểm có thể tồn tại ở dạng có dấu, không dấu, viết tắt hoặc pha trộn tiếng Việt và tiếng Anh \cite{yeboah2021}. Địa chỉ thường không hoàn toàn chuẩn hóa theo cùng một cấu trúc, đặc biệt khi dữ liệu đến từ cộng đồng, bản đồ mở hoặc nhiều nhà cung cấp khác nhau \cite{yeboah2021}. Danh mục địa điểm cũng có thể không đồng nhất, ví dụ cùng một điểm được gắn là quán cà phê, nhà hàng, điểm tham quan hoặc dịch vụ du lịch tùy nguồn dữ liệu \cite{zhang2023}. Ngoài ra, dữ liệu mở thường có mức độ đầy đủ khác nhau theo vùng, nên hệ thống cần lưu nguồn gốc và đánh giá chất lượng theo từng bản ghi thay vì giả định mọi khu vực đều có dữ liệu đồng đều \cite{yeboah2021}.\par

\subsection{Quy trình xử lý và kiến trúc nhiều lớp}

Quy trình trích xuất, biến đổi và nạp dữ liệu (ETL) được dùng làm khung xử lý chính cho dự án \cite{ibmETL}. Bước trích xuất lấy dữ liệu từ nguồn bản đồ, API hoặc dữ liệu nội bộ; bước biến đổi làm sạch, chuẩn hóa, kiểm tra và hợp nhất dữ liệu; bước nạp đưa dữ liệu đã xử lý vào hệ thống lưu trữ hoặc lớp phục vụ \cite{ibmETL}. Dự án ưu tiên xử lý theo lô vì dữ liệu địa điểm du lịch thường không yêu cầu cập nhật tức thời như dữ liệu giao dịch hoặc cảm biến \cite{awsStreaming}.\par

Kiến trúc nhiều lớp giúp tách dữ liệu theo mức độ tin cậy: dữ liệu thô, dữ liệu đã chuẩn hóa và dữ liệu sẵn sàng phục vụ \cite{databricksMedallion}. Cách tổ chức này bảo toàn dữ liệu nguồn, hỗ trợ tái xử lý, giảm rủi ro đưa dữ liệu lỗi lên giao diện và giúp truy ngược khi dữ liệu đầu ra có vấn đề \cite{ibmLineage}.\par

\subsection{Chất lượng, quản trị và quan sát dữ liệu}

Chất lượng dữ liệu quyết định khả năng sử dụng dữ liệu trong phân tích và ra quyết định \cite{ibmDataQuality}. Với dữ liệu POI, các lỗi thường gặp gồm thiếu tọa độ, thiếu tên, sai danh mục, trùng lặp, thiếu thông tin liên hệ hoặc khác biệt mô tả giữa các nguồn \cite{yeboah2021}. Do đó, dự án sử dụng cơ chế chấm điểm chất lượng và vùng chờ xem xét để kết hợp tự động hóa với kiểm soát thủ công khi cần \cite{ibmDataQuality}.\par

Truy vết dữ liệu giúp hệ thống biết một bản ghi đến từ nguồn nào, được thu thập khi nào, đã đi qua quy tắc xử lý nào và có liên hệ với bản ghi ở các lớp trước ra sao \cite{ibmLineage}. Quản trị dữ liệu bổ sung các nguyên tắc về trách nhiệm, chính sách sử dụng, metadata, bảo mật và khả năng duy trì dữ liệu trong toàn bộ vòng đời hệ thống \cite{damaDmbok}. Quan sát dữ liệu mở rộng giám sát kỹ thuật sang các chỉ số như độ tươi mới, tỷ lệ bản ghi đạt chất lượng, lỗi theo nguồn và trạng thái pipeline \cite{opentelemetrySignals}.\par


\section{Cơ sở công nghệ}

Cơ sở công nghệ của dự án được lựa chọn theo nguyên tắc vừa đủ: đáp ứng yêu cầu tích hợp và phục vụ dữ liệu POI trong phạm vi hiện tại, nhưng vẫn giữ khả năng mở rộng khi quy mô dữ liệu hoặc yêu cầu vận hành tăng lên \cite{ibmWarehouseLake}. Nhóm công nghệ chính gồm MongoDB cho lưu trữ document, Python/FastAPI cho xử lý dữ liệu, Node.js/Express cho API trung gian, React cho dashboard và Docker Compose cho đóng gói môi trường \cite{dockerCompose}.\par

\subsection{Lưu trữ và xử lý dữ liệu phía máy chủ}

MongoDB phù hợp với dữ liệu POI vì mô hình document hỗ trợ dữ liệu bán cấu trúc, nhiều trường tùy chọn và dữ liệu JSON từ các API bên ngoài \cite{mongodbDataModel}. Cách lưu trữ này giúp dự án tiếp nhận dữ liệu thô linh hoạt ở giai đoạn đầu, sau đó chuẩn hóa dần qua các lớp xử lý trước khi phục vụ dashboard hoặc API \cite{mongodbDataModel}. Aggregation pipeline của MongoDB đáp ứng các nhu cầu thống kê hiện tại như đếm theo thành phố, danh mục, nguồn dữ liệu hoặc mức chất lượng \cite{mongodbAggregation}.\par

Python được dùng cho các tác vụ xử lý dữ liệu, tích hợp API và xây dựng pipeline \cite{pythonApps}. FastAPI hỗ trợ xây dựng service gọn, có tài liệu API tự động và phù hợp với type hint của Python \cite{fastapiFeatures}. Pydantic giúp kiểm tra dữ liệu đầu vào, cấu hình và dữ liệu trả về ở ranh giới dịch vụ \cite{pydanticModels}.\par

\subsection{API, giao diện và hợp đồng dữ liệu}

Node.js, Express và TypeScript được dùng cho lớp API trung gian để tách frontend khỏi dữ liệu thô và gom các logic phục vụ dashboard \cite{nodeAbout}. OpenAPI được dùng để mô tả hợp đồng API, bao gồm endpoint, tham số, request, response và schema dữ liệu trao đổi giữa backend và frontend \cite{openapiSpec}. Zod hỗ trợ kiểm tra dữ liệu runtime ở phía TypeScript, giúp phát hiện dữ liệu không đúng cấu trúc trước khi đi sâu vào logic giao diện hoặc API \cite{zodDocs}. Orval hỗ trợ sinh mã client từ đặc tả OpenAPI, qua đó giảm thao tác viết tay và giảm nguy cơ lệch kiểu dữ liệu giữa backend và frontend \cite{orvalDocs}.\par

React, Vite, TanStack Query, Tailwind CSS, Radix UI và Recharts tạo thành nhóm công nghệ phù hợp cho dashboard dạng ứng dụng một trang \cite{reactDocs}. Nhóm này hỗ trợ xây dựng giao diện theo thành phần, quản lý trạng thái dữ liệu server, tạo bảng/bộ lọc/biểu đồ và trình bày các KPI vận hành dữ liệu \cite{tableauKpi}.\par

\subsection{Đóng gói và hướng mở rộng}

Docker Compose được dùng để chạy nhiều thành phần như frontend, API server, service ETL và cơ sở dữ liệu trong cùng một môi trường thống nhất \cite{dockerCompose}. Cách này giảm khác biệt giữa môi trường phát triển, trình diễn và tài liệu hướng dẫn \cite{dockerMultiContainer}.\par

Các công nghệ như Spark, Airflow, Kafka, Iceberg, Trino hoặc Kubernetes chưa được đưa vào triển khai chính vì phạm vi hiện tại chưa cần hạ tầng phân tán phức tạp \cite{sparkDocs}. Khi dữ liệu tăng mạnh hoặc hệ thống cần xử lý gần thời gian thực, các công nghệ này có thể được bổ sung để xử lý dữ liệu lớn, điều phối DAG phức tạp hoặc mở rộng kiến trúc lakehouse \cite{airflowDocs}.\par

\section{Lý do lựa chọn và định hướng áp dụng}

Các nghiên cứu về hệ thống gợi ý du lịch thường nhấn mạnh thuật toán đề xuất hơn là nền tảng dữ liệu đầu vào \cite{zhang2023}. Trong thực tế, thuật toán hoặc dashboard khó tạo ra kết quả đáng tin nếu POI bị trùng, thiếu tọa độ, sai danh mục hoặc không rõ nguồn gốc \cite{ibmDataQuality}. Vì vậy, dự án đặt nền tảng dữ liệu POI có kiểm soát chất lượng, truy vết và tích hợp nhiều nguồn làm trọng tâm trước khi phát triển các chức năng phân tích hoặc gợi ý nâng cao \cite{zhang2023}.\par

\subsection{So sánh với cách xử lý thủ công và lưu trữ trực tiếp}

Cách làm đơn giản nhất đối với dữ liệu POI là thu thập dữ liệu, lưu vào một bảng hoặc một file trung gian, sau đó cho dashboard đọc trực tiếp \cite{ibmETL}. Ưu điểm của cách này là dễ triển khai, ít thành phần kỹ thuật, phù hợp với kiểm thử nhanh hoặc trình diễn ban đầu \cite{ibmETL}. Tuy nhiên, khi dữ liệu đến từ nhiều nguồn, cách lưu trực tiếp làm tăng nguy cơ trộn lẫn dữ liệu thô với dữ liệu đã xử lý, khiến hệ thống khó biết bản ghi nào đã được kiểm tra và bản ghi nào còn thiếu thông tin \cite{ibmDataQuality}. Nếu không có lớp trung gian để chuẩn hóa và kiểm tra, các lỗi như trùng địa điểm, sai danh mục, thiếu tọa độ hoặc sai nguồn có thể đi thẳng lên dashboard \cite{yeboah2021}.\par

Một phương án cũ khác là xử lý thủ công bằng file bảng tính hoặc script rời rạc \cite{ibmETL}. Ưu điểm của hướng này là linh hoạt, dễ sửa nhanh và không cần thiết kế hạ tầng phức tạp \cite{ibmETL}. Nhược điểm là khó tái lập kết quả, khó kiểm soát phiên bản dữ liệu, khó xác định quy tắc nào đã thay đổi bản ghi và khó mở rộng khi số nguồn hoặc số thành phố tăng lên \cite{ibmLineage}. Vì vậy, dự án không chọn cách xử lý thủ công làm hướng chính, mà chuyển các bước thu thập, chuẩn hóa, chấm điểm và phục vụ dữ liệu thành pipeline có cấu trúc \cite{ibmETL}.\par

\subsection{So sánh với cơ sở dữ liệu quan hệ và data warehouse}

Cơ sở dữ liệu quan hệ là lựa chọn quen thuộc vì hỗ trợ ràng buộc dữ liệu, truy vấn có cấu trúc và mô hình bảng rõ ràng \cite{ibmWarehouseLake}. Với dữ liệu nghiệp vụ ổn định, schema quan hệ giúp kiểm soát kiểu dữ liệu và quan hệ giữa các bảng tốt \cite{ibmWarehouseLake}. Tuy nhiên, dữ liệu POI thường có nhiều trường tùy chọn như giờ mở cửa, ảnh, số điện thoại, website, đánh giá, loại hình dịch vụ hoặc metadata theo nguồn \cite{mongodbDataModel}. Nếu ép dữ liệu này vào schema quan hệ chặt ngay từ đầu, dự án phải tạo nhiều bảng phụ, nhiều quan hệ và nhiều lần migration khi nguồn dữ liệu thay đổi \cite{mongodbDataModel}.\par

MongoDB được chọn vì phù hợp hơn với giai đoạn tích hợp dữ liệu bán cấu trúc \cite{mongodbDataModel}. Dữ liệu thô từ API có thể được lưu gần với dạng ban đầu, sau đó chuẩn hóa dần qua các bước xử lý và đưa bản ghi đủ chất lượng vào lớp phục vụ \cite{mongodbDataModel}. Lựa chọn này không phủ nhận giá trị của cơ sở dữ liệu quan hệ, mà phản ánh đặc điểm hiện tại của bài toán: schema POI còn thay đổi, dữ liệu đến từ nhiều nguồn và cần ưu tiên khả năng tiếp nhận linh hoạt \cite{yeboah2021}.\par

Data warehouse truyền thống mạnh trong phân tích dữ liệu có cấu trúc, báo cáo định kỳ và các mô hình dữ liệu đã ổn định \cite{ibmWarehouseLake}. Ưu điểm của hướng này là khả năng tổ chức dữ liệu phục vụ phân tích tốt, phù hợp với doanh nghiệp có quy trình dữ liệu trưởng thành và yêu cầu báo cáo rõ ràng \cite{ibmWarehouseLake}. Tuy nhiên, với dự án quy mô nhỏ đến trung bình, triển khai warehouse đầy đủ có thể làm tăng chi phí thiết kế mô hình, vận hành và nạp dữ liệu trước khi chứng minh được giá trị của pipeline POI \cite{ibmWarehouseLake}. Vì vậy, dự án chỉ kế thừa tư duy phân lớp dữ liệu và chất lượng dữ liệu, chưa triển khai một data warehouse hoàn chỉnh \cite{databricksMedallion}.\par

\subsection{So sánh với streaming, lakehouse và nền tảng dữ liệu lớn}

Streaming phù hợp với hệ thống cần phản ứng gần thời gian thực, ví dụ sự kiện người dùng, log liên tục, cảm biến hoặc dữ liệu giao dịch tốc độ cao \cite{awsStreaming}. Ưu điểm của streaming là giảm độ trễ và cho phép hệ thống xử lý sự kiện ngay khi phát sinh \cite{awsStreaming}. Tuy nhiên, streaming cũng kéo theo độ phức tạp về message broker, checkpoint, retry, xử lý thứ tự sự kiện, giám sát độ trễ và cơ chế phục hồi khi lỗi \cite{awsStreaming}. Với dữ liệu POI du lịch, nhiều thông tin như tên địa điểm, danh mục, địa chỉ hoặc website thường không thay đổi từng giây, nên batch processing hợp lý hơn cho phạm vi hiện tại \cite{ibmETL}.\par

Lakehouse là hướng hiện đại để kết hợp ưu điểm của data lake và data warehouse trên quy mô lớn \cite{armbrust2021}. Ưu điểm của lakehouse là hỗ trợ dữ liệu lớn, nhiều định dạng, phân tích nâng cao và khả năng mở rộng tốt \cite{armbrust2021}. Tuy nhiên, nếu áp dụng đầy đủ lakehouse ngay từ đầu, dự án phải xử lý thêm nhiều vấn đề về object storage, định dạng bảng, engine truy vấn, orchestration và vận hành hạ tầng \cite{armbrust2021}. Với mục tiêu hiện tại là chứng minh pipeline POI có kiểm soát chất lượng và dashboard vận hành, kiến trúc gọn nhẹ bằng MongoDB, FastAPI, Express, React và Docker Compose phù hợp hơn \cite{dockerCompose}.\par

Các công nghệ như Spark, Airflow, Kafka, Iceberg, Trino hoặc Kubernetes vẫn có giá trị cho giai đoạn mở rộng \cite{sparkDocs}. Spark phù hợp khi khối lượng dữ liệu lớn hơn khả năng xử lý của một service thông thường \cite{sparkDocs}. Airflow phù hợp khi pipeline có nhiều DAG, nhiều phụ thuộc và cần lịch chạy phức tạp \cite{airflowDocs}. Kafka phù hợp khi hệ thống cần event streaming hoặc tích hợp nhiều producer và consumer thời gian thực \cite{awsStreaming}. Tuy nhiên, đưa toàn bộ các công nghệ này vào phiên bản hiện tại sẽ làm tăng độ khó cài đặt, giải thích, kiểm thử và bảo trì mà chưa tạo ra lợi ích tương xứng với quy mô dữ liệu \cite{dockerCompose}.\par

\subsection{Lý do chọn hướng kiến trúc vừa đủ}

Dự án chọn hướng tiếp cận vừa đủ thay vì triển khai ngay một nền tảng dữ liệu lớn hoàn chỉnh \cite{armbrust2021}. Lựa chọn này phù hợp với mục tiêu xây dựng được pipeline dữ liệu du lịch có kiểm soát, có dashboard vận hành, có API phục vụ và có khả năng giải thích trong phạm vi dự án \cite{ibmETL}. Hướng này giữ lại các nguyên tắc quan trọng của hệ thống dữ liệu hiện đại như phân lớp dữ liệu, kiểm soát chất lượng, truy vết, hợp đồng API và quan sát vận hành, nhưng triển khai bằng các công nghệ nhẹ hơn để giảm chi phí vận hành \cite{damaDmbok}.\par

Ưu điểm lớn nhất của hướng vừa đủ là cân bằng giữa tính học thuật, tính kỹ thuật và tính khả thi \cite{ibmWarehouseLake}. Về học thuật, hệ thống vẫn bám vào các khái niệm nền tảng như ETL, chất lượng dữ liệu, lineage, governance và API contract \cite{damaDmbok}. Về kỹ thuật, hệ thống có đủ thành phần để thu thập, xử lý, lưu trữ, phục vụ và quan sát dữ liệu POI \cite{ibmETL}. Về tính khả thi, nhóm phát triển có thể triển khai, kiểm thử và trình diễn hệ thống mà không phải vận hành một cụm dữ liệu phân tán phức tạp \cite{dockerCompose}.\par

So với các phương pháp cũ, hướng này giảm rủi ro xử lý dữ liệu cảm tính và giảm phụ thuộc vào thao tác thủ công \cite{ibmDataQuality}. So với nền tảng dữ liệu lớn, hướng này tránh được việc triển khai quá mức trong khi vẫn giữ đường mở rộng rõ ràng \cite{armbrust2021}. Khi dữ liệu tăng hoặc yêu cầu vận hành thay đổi, hệ thống có thể bổ sung Airflow để điều phối pipeline, Kafka để xử lý sự kiện, Spark để xử lý dữ liệu lớn hoặc lakehouse để lưu trữ phân tích quy mô lớn \cite{airflowDocs}.\par

\subsection{Định hướng áp dụng trong dự án}

Định hướng áp dụng của dự án là xây dựng một nền tảng dữ liệu POI có thể vận hành được trong bối cảnh dữ liệu du lịch Việt Nam, bao gồm tên địa điểm có dấu/không dấu, danh mục đa dạng, địa chỉ không đồng nhất và khác biệt chất lượng giữa các nguồn \cite{yeboah2021}. Pipeline cần ưu tiên khả năng tiếp nhận nhiều nguồn, lưu nguồn gốc bản ghi, chuẩn hóa từng trường quan trọng, đánh giá chất lượng và chỉ đưa dữ liệu đủ tin cậy vào lớp phục vụ \cite{ibmDataQuality}. Dashboard không chỉ hiển thị số lượng POI mà còn cần thể hiện trạng thái pipeline, tỷ lệ bản ghi đạt chất lượng, nguồn dữ liệu và các bất thường cần kiểm tra \cite{tableauKpi}.\par

\subsubsection{Tiêu chí Gold, Pending Review và Quarantine}

Bản ghi được đưa vào Gold khi có đủ các trường cốt lõi như tên địa điểm, tọa độ hợp lệ, khu vực hoặc thành phố, danh mục đã chuẩn hóa và nguồn dữ liệu rõ ràng \cite{ibmDataQuality}. Bản ghi Gold cũng cần không trùng lặp hoặc đã được hợp nhất với các bản ghi tương ứng từ nguồn khác, để dashboard và API có thể dùng như dữ liệu phục vụ chính \cite{ibmLineage}. Trạng thái Gold phù hợp với các POI đạt ngưỡng chất lượng đã định và đủ tin cậy để xuất hiện trong thống kê, bộ lọc, bản đồ hoặc endpoint công khai của hệ thống \cite{tableauKpi}.\par

Bản ghi được đưa vào Pending Review khi có giá trị sử dụng nhưng còn thiếu hoặc chưa chắc chắn ở một số trường quan trọng, chẳng hạn thiếu địa chỉ chi tiết, thiếu thông tin liên hệ, danh mục chưa rõ hoặc tọa độ cần kiểm tra thêm \cite{ibmDataQuality}. Nhóm Pending Review giúp hệ thống không loại bỏ quá sớm các bản ghi có tiềm năng, đồng thời tạo vùng kiểm soát cho việc bổ sung dữ liệu, kiểm tra thủ công hoặc xử lý lại ở lần chạy sau \cite{ibmLineage}.\par

Bản ghi được đưa vào Quarantine khi có lỗi nghiêm trọng như thiếu tên, thiếu tọa độ, tọa độ nằm ngoài phạm vi hợp lý, nguồn dữ liệu không xác định hoặc xung đột trùng lặp chưa thể hợp nhất \cite{ibmDataQuality}. Dữ liệu trong Quarantine không được dùng trực tiếp cho dashboard hoặc API phục vụ người dùng cho đến khi được sửa lỗi, bổ sung nguồn gốc hoặc xử lý lại bằng quy tắc phù hợp \cite{openapiSpec}.\par

Trong giai đoạn hiện tại, hệ thống tập trung vào quy trình batch có kiểm soát, lưu trữ document linh hoạt, API rõ hợp đồng và dashboard vận hành \cite{ibmETL}. Trong giai đoạn tiếp theo, hệ thống có thể mở rộng theo ba hướng: thêm nguồn dữ liệu mới, mở rộng địa bàn thu thập và bổ sung năng lực phân tích nâng cao \cite{zhang2023}. Khi các hướng này phát triển, nền tảng hiện tại vẫn giữ vai trò định hướng: dữ liệu cần được thu thập có kiểm soát, xử lý có truy vết, đánh giá có tiêu chí và phục vụ qua hợp đồng rõ ràng \cite{openapiSpec}.\par
\renewcommand{\refname}{Tài liệu tham khảo}
\begin{thebibliography}{99}
\addcontentsline{toc}{section}{Tài liệu tham khảo}
\item[]\textbf{Nhóm tài liệu lý thuyết, nghiên cứu và tiêu chuẩn}

\bibitem{armbrust2021} Armbrust, M., Ghodsi, A., Xin, R., \& Zaharia, M. (2021). \textit{Lakehouse: A new generation of open platforms that unify data warehousing and advanced analytics}. CIDR. \url{https://www.databricks.com/sites/default/files/2020/12/cidr_lakehouse.pdf}
\bibitem{damaDmbok} DAMA International. (n.d.). \textit{DAMA-DMBOK data management body of knowledge}. Retrieved May 18, 2026, from \url{https://dama.org/learning-resources/dama-data-management-body-of-knowledge-dmbok/}
\bibitem{ibmWarehouseLake} IBM. (n.d.). \textit{Data warehouses vs. data lakes vs. data lakehouses}. Retrieved May 18, 2026, from \url{https://www.ibm.com/think/topics/data-warehouse-vs-data-lake-vs-data-lakehouse}
\bibitem{ibmDataQuality} IBM. (n.d.). \textit{What are data quality dimensions?} Retrieved May 18, 2026, from \url{https://www.ibm.com/think/topics/data-quality-dimensions}
\bibitem{ibmLineage} IBM. (n.d.). \textit{What is data lineage?} Retrieved May 18, 2026, from \url{https://www.ibm.com/topics/data-lineage}
\bibitem{ibmETL} IBM. (n.d.). \textit{What is ETL (extract, transform, load)?} Retrieved May 18, 2026, from \url{https://www.ibm.com/topics/etl}
\bibitem{openapiSpec} OpenAPI Initiative. (n.d.). \textit{The OpenAPI Specification explained}. Retrieved May 18, 2026, from \url{https://learn.openapis.org/specification/}
\bibitem{opentelemetrySignals} OpenTelemetry. (n.d.). \textit{Signals}. Retrieved May 18, 2026, from \url{https://opentelemetry.io/docs/concepts/signals/}
\bibitem{tableauKpi} Tableau. (n.d.). \textit{What is a KPI dashboard?} Retrieved May 18, 2026, from \url{https://www.tableau.com/kpi/what-is-kpi-dashboard}
\bibitem{yeboah2021} Yeboah, G., Porto de Albuquerque, J., Troilo, R., Tregonning, G., Perera, S., \& Ahmed, S. A. K. S. (2021). Analysis of OpenStreetMap data quality at different stages of a participatory mapping process. \textit{ISPRS International Journal of Geo-Information, 10}(4), Article 265. \url{https://www.mdpi.com/2220-9964/10/4/265}
\bibitem{zhang2023} Zhang, J., Chow, C.-Y., \& Li, Y. (2023). \textit{A survey on point-of-interest recommendations leveraging heterogeneous data}. arXiv. \url{https://arxiv.org/abs/2308.07426}
\item[]\textbf{Nhóm tài liệu công nghệ và công cụ}

\bibitem{awsStreaming} Amazon Web Services. (n.d.). \textit{What is streaming data?} Retrieved May 18, 2026, from \url{https://aws.amazon.com/streaming-data/}
\bibitem{airflowDocs} Apache Airflow. (n.d.). \textit{Apache Airflow documentation}. Retrieved May 18, 2026, from \url{https://airflow.apache.org/docs/}
\bibitem{sparkDocs} Apache Spark. (n.d.). \textit{Apache Spark overview}. Retrieved May 18, 2026, from \url{https://spark.apache.org/docs/latest/}
\bibitem{databricksMedallion} Databricks. (n.d.). \textit{What is the medallion lakehouse architecture?} Retrieved May 18, 2026, from \url{https://docs.databricks.com/en/lakehouse/medallion.html}
\bibitem{dockerCompose} Docker. (n.d.). \textit{Docker Compose}. Retrieved May 18, 2026, from \url{https://docs.docker.com/compose/}
\bibitem{dockerMultiContainer} Docker. (n.d.). \textit{Multi-container applications}. Retrieved May 18, 2026, from \url{https://docs.docker.com/get-started/docker-concepts/running-containers/multi-container-applications/}
\bibitem{fastapiFeatures} FastAPI. (n.d.). \textit{Features}. Retrieved May 18, 2026, from \url{https://fastapi.tiangolo.com/features/}
\bibitem{mongodbAggregation} MongoDB. (n.d.). \textit{Aggregation pipeline}. Retrieved May 18, 2026, from \url{https://www.mongodb.com/docs/manual/core/aggregation-pipeline/}
\bibitem{mongodbDataModel} MongoDB. (n.d.). \textit{Data modeling in MongoDB}. Retrieved May 18, 2026, from \url{https://www.mongodb.com/docs/master/core/data-model-design/}
\bibitem{nodeAbout} OpenJS Foundation. (n.d.). \textit{About Node.js}. Retrieved May 18, 2026, from \url{https://nodejs.org/en/about}
\bibitem{orvalDocs} Orval. (n.d.). \textit{Overview}. Retrieved May 18, 2026, from \url{https://orval.dev/overview}
\bibitem{pydanticModels} Pydantic. (n.d.). \textit{Models}. Retrieved May 18, 2026, from \url{https://docs.pydantic.dev/latest/concepts/models/}
\bibitem{pythonApps} Python Software Foundation. (n.d.). \textit{Applications for Python}. Retrieved May 18, 2026, from \url{https://www.python.org/about/apps/}
\bibitem{reactDocs} React. (n.d.). \textit{React documentation}. Retrieved May 18, 2026, from \url{https://react.dev/}
\bibitem{zodDocs} Zod. (n.d.). \textit{Introduction}. Retrieved May 18, 2026, from \url{https://zod.dev/}

\end{thebibliography}

\end{document}
